\section{Attachment cohomology and the separation ledger}
\label{sec:separation-budget}

The realization theorem uses a finite tree meeting each atom at one point.
There are two logically distinct requirements hidden in that choice: the
connector itself has no holes, and every attachment trace is connected.  The
next theorem separates these requirements and gives an exact attachment
ledger.  An extra component of a trace creates a complementary component,
whereas a complementary component of the trace itself fills one separation
already carried by the connector.  Thus attachments can create and destroy
holes, but only through two explicitly measured trace defects.

For two nonseparating plane continua, the connected-intersection
sufficiency is Janiszewski's theorem and the converse belongs to the
classical Janiszewski--Mullikin theory \cite{Jan15,Mul22,New51,Why42}.  The
form needed here keeps arbitrary initial connector cohomology, allows
separating traces, records the homomorphism killed by each attachment, and
iterates over a disjoint atom family.  It also sits naturally in the
classical unicoherence and multicoherence tradition \cite{Bor31,Eil36}:
the term $c_i-1$ is the classical multicoherence contribution of the
decomposition $B_i=B_{i-1}\cup A_i$.  The contribution here is its
simultaneous Mayer--Vietoris bookkeeping and carrier application; we make
no priority claim for the underlying contact-count invariant.

For a proper compact set $K\subset\Chat$ whose complement has finitely many
components, define its \emph{separation excess} by
\begin{equation}
 \sigma(K):=\#\pi_0(\Chat\setminus K)-1.
 \label{eq:G-separation-excess}
\end{equation}
Thus $\sigma(K)=0$ exactly when $K$ is nonseparating.  Since
$\Chat\setminus K$ is open in a surface, it is locally path connected, so its
connected components and path components agree.  Alexander duality for
compacta therefore identifies
\begin{equation}
 \sigma(K)=\operatorname{rank}\check H^1(K;\Z).
 \label{eq:G-dual-rank}
\end{equation}
We use the compact-pair Mayer--Vietoris sequence for \v{C}ech cohomology; both
facts are recalled in the generality needed here in \cite[Chapter~6]{Spa66}.

\begin{theorem}[Attachment exact sequence and separation ledger]
\label{thm:separation-budget}
Let $s\ge1$, and let $C,A_1,\dots,A_s\subset\C$ be compact connected sets.
Assume that the
atoms $A_i$ are pairwise disjoint full continua and that
\[
 T_i:=C\cap A_i
\]
is nonempty.  Put $H=C\cup\bigcup_iA_i$, and let
$B_0=C$ and $B_i=C\cup A_1\cup\cdots\cup A_i$.  For each $i$, the
restriction homomorphism
\[
 r_i\colon\check H^1(B_{i-1};\Z)\longrightarrow
 \check H^1(T_i;\Z)
\]
is surjective, and there is a natural short exact
sequence
\begin{equation}
 0\longrightarrow
 \widetilde{\check H}^{\,0}(T_i;\Z)
 \xrightarrow{\ \delta_i\ }
 \check H^1(B_i;\Z)
 \longrightarrow
 \ker r_i
 \longrightarrow0.
 \label{eq:G-attachment-exact}
\end{equation}
If $T_i$ is nonseparating for every $i$, this specializes, noncanonically,
to
\begin{equation}
 \check H^1(H;\Z)
 \cong
 \check H^1(C;\Z)\oplus
 \bigoplus_{i=1}^{s}
 \widetilde{\check H}^{\,0}(T_i;\Z).
 \label{eq:G-attachment-splitting}
\end{equation}

Suppose now that $\Chat\setminus C$ and every
$\Chat\setminus T_i$ have finitely many components, and that $T_i$ has
finitely many connected components, say $c_i$.  Then all groups in
\eqref{eq:G-attachment-exact} have finite rank, all complements
$\Chat\setminus B_i$ have finitely many components, and
\begin{equation}
 \boxed{\quad
 \sigma(H)=\sigma(C)+
 \sum_{i=1}^{s}\bigl(c_i-1-\sigma(T_i)\bigr).
 \quad}
 \label{eq:G-budget}
\end{equation}
In particular, at the $i$-th attachment
\begin{equation}
 \sigma(T_i)\le \sigma(B_{i-1}),
 \qquad
 \sigma(B_i)=\sigma(B_{i-1})+c_i-1-\sigma(T_i).
 \label{eq:G-step-ledger}
\end{equation}
\end{theorem}

\begin{proof}
Every $B_i$ is connected.  Indeed, $B_0=C$ is connected, and at the
$i$-th step the connected set $A_i$ meets $B_{i-1}$ in the nonempty set
$T_i$.  Pairwise disjointness of the atoms gives the exact intersection
identity
\[
 B_{i-1}\cap A_i
 =
 (C\cap A_i)\cup\bigcup_{j<i}(A_j\cap A_i)
 =T_i .
\]

Apply the reduced compact-pair Mayer--Vietoris sequence to the closed cover
$B_i=B_{i-1}\cup A_i$.  Its relevant segment is
\[
\begin{split}
 \widetilde{\check H}^{\,0}(B_{i-1})
 \oplus\widetilde{\check H}^{\,0}(A_i)
 &\longrightarrow
 \widetilde{\check H}^{\,0}(T_i)
 \xrightarrow{\delta_i}
 \check H^1(B_i)\\
 &\longrightarrow
 \check H^1(B_{i-1})\oplus\check H^1(A_i)
 \longrightarrow
 \check H^1(T_i).
\end{split}
\]
The first two reduced groups vanish because $B_{i-1}$ and $A_i$ are
connected.  Fullness of $A_i$ and Alexander duality give
$\check H^1(A_i;\Z)=0$.  The next term of the sequence is
$\check H^2(B_i;\Z)$, and it vanishes: $B_i$ is a compact subset of the
plane, hence a proper compact subset of $S^2$, and compacta Alexander
duality identifies its reduced second cohomology with reduced homology in
degree $-1$ of the nonempty complement.  Exactness therefore makes $r_i$
surjective and leaves exactly \eqref{eq:G-attachment-exact}.

Alexander duality identifies $\check H^1(B_{i-1};\Z)$ with the reduced
zeroth homology of the open surface $\Chat\setminus B_{i-1}$, hence with a
free abelian group.  Its subgroup $\ker r_i$ is free and therefore
projective, so \eqref{eq:G-attachment-exact} splits, although not
canonically.  When every $T_i$ is nonseparating, the target of $r_i$
vanishes and iteration gives \eqref{eq:G-attachment-splitting}.

Under the finiteness hypotheses, the components of $T_i$ are clopen and
\[
 \operatorname{rank}\widetilde{\check H}^{\,0}(T_i;\Z)=c_i-1,
 \qquad
 \operatorname{rank}\check H^1(T_i;\Z)=\sigma(T_i).
\]
Surjectivity of $r_i$ gives the first inequality in
\eqref{eq:G-step-ledger}; taking ranks in
\eqref{eq:G-attachment-exact} gives its equality.  Starting from the
finite-rank group $\check H^1(C;\Z)$, induction makes every
$\check H^1(B_i;\Z)$ finite rank.  Alexander duality then makes every
complement $\Chat\setminus B_i$ have finitely many components.  Summing
the step equalities gives \eqref{eq:G-budget}.
\end{proof}

The ledger has a useful geometric reading.  For example, let $C$ be the
unit circle together with the radial segment $[1,2]$, and attach the full
closed unit disk $A$.  The trace $C\cap A$ is the unit circle.  It is
connected but has separation excess one, so the attachment fills the one
bounded complementary component of $C$:
\[
 \sigma(C\cup A)=1+(1-1)-1=0.
\]
In contrast, a nonseparating trace cannot cancel connector homology; its
only possible contribution is the creation term $c_i-1$.

The trace-separation term vanishes automatically for the carrier geometry
that matters most here.

\begin{lemma}[Intersections of full compacta are nonseparating]
\label{lem:full-intersection}
If $K,L\subset\C$ are compact and full, then $K\cap L$ is nonseparating
(whether or not it is connected).  In particular
$\check H^1(K\cap L;\Z)=0$.
\end{lemma}

\begin{proof}
If $K\cap L=\varnothing$, its complement is $\C$ and
$\check H^1(\varnothing;\Z)=0$, so the conclusion is immediate.  Assume
$K\cap L\ne\varnothing$.
The two connected sets $\C\setminus K$ and $\C\setminus L$ contain a common
neighborhood of infinity.  Their union is therefore connected, and
\[
 \C\setminus(K\cap L)
 =(\C\setminus K)\cup(\C\setminus L)
\]
is connected.  Alexander duality gives the cohomological conclusion.
\end{proof}

\begin{corollary}[Exact criterion for a full continuum connector]
\label{cor:full-connector-traces}
In Theorem~\ref{thm:separation-budget}, suppose that $C$ is also a full
continuum.  Then every trace is nonseparating, and $H$ is full if and only
if every attachment trace $C\cap A_i$ is connected.  If, in addition,
$T_i$ has finitely many connected components $c_i$ for every $i$, then
\begin{equation}
 \#\pi_0(\Chat\setminus H)
 =
 1+\sum_{i=1}^{s}(c_i-1).
 \label{eq:G-full-budget}
\end{equation}
\end{corollary}

\begin{proof}
Lemma~\ref{lem:full-intersection}, applied to $C$ and $A_i$, verifies
that $\check H^1(T_i;\Z)=0$.  The splitting
\eqref{eq:G-attachment-splitting}, together with
$\check H^1(C;\Z)=0$, gives
\[
 \check H^1(H;\Z)
 \cong
 \bigoplus_i\widetilde{\check H}^{\,0}(T_i;\Z).
\]
For a compact Hausdorff space, reduced \v{C}ech cohomology in degree zero
vanishes exactly when the space is connected.  Alexander duality now gives
the stated if-and-only-if criterion without a finiteness assumption.
When all $c_i$ are finite, every $T_i$ has $\sigma(T_i)=0$ and
$\sigma(C)=0$, so \eqref{eq:G-budget} becomes
\eqref{eq:G-full-budget}.
\end{proof}

Corollary~\ref{cor:full-connector-traces} both strengthens and sharpens
the one-point tree formulation of Lemma~\ref{lem:tree-full}, and its
if-and-only-if clause recovers the Janiszewski--Mullikin criterion in this
finite-trace setting.  The connector need not be a tree, locally
connected, or one dimensional: any full continuum is allowed.  Singleton
traces can be replaced by arbitrary connected traces.  Conversely,
disconnected traces are not merely inconvenient.  When the connector is
full they force holes, and the number forced is exact.

\begin{corollary}[Protected realization by a continuum connector]
\label{cor:continuum-connector}
Retain the Jordan localization, atoms, obstacle set, and protected
crosscut chains of Section~\ref{sec:certificates}, with pairwise disjoint
relative terminal collar tails.  Let $C\Subset U$ be a full continuum such
that:
\begin{enumerate}
\item $C\cap Z=\varnothing$;
\item $C\cap A_i$ is nonempty and connected for every $i$; and
\item
\[
 (C\setminus A_i)\cap
 \bigl(\overline{D_{i,a,1}}\setminus A_i\bigr)=\varnothing
 \quad\text{for every protected chain }(i,a).
\]
\end{enumerate}
Then
\[
 H_C:=C\cup\bigcup_iA_i
\]
is a compact connected full carrier avoiding $Z$.  Every protected chain
survives in $\Chat\setminus H_C$; the survival assignment is injective on
families with disjoint relative terminal collar tails and preserves the
cyclic order carried by each common source complement.  No local
connectivity of $C$ or of any atom is required.
\end{corollary}

\begin{proof}
Corollary~\ref{cor:full-connector-traces} makes $H_C$ full, and the nonempty
traces make it connected.  The collar-clearance hypothesis for each chain is
exactly the hypothesis of Lemma~\ref{lem:crosscut-survival}, because the
other atoms miss its relative terminal collar tail.  That lemma supplies
survival, injectivity, and cyclic-order preservation.  Compact containment
and obstacle avoidance follow from the stated hypotheses and the finite
union.
\end{proof}

\subsection{Sharpness and defect birth at a limit}
\label{subsec:G-limit-defect}

The separation ledger exposes a compactness obstruction that is invisible at
every finite stage.  Even a sequence of tree carriers with singleton traces
can converge to a carrier with arbitrarily many bounded complementary
components.

\begin{proposition}[Arbitrary separation defect born in a Hausdorff limit]
\label{prop:trace-birth}
For every integer $m\ge1$ there are a full atom $A$, full arc connectors
$C_n$, and a full arc $C$ such that
\[
 C_n\longrightarrow C
 \quad\text{and}\quad
 A\cup C_n\longrightarrow A\cup C
\]
in the Hausdorff metric,
\[
 A\cap C_n\ \text{is a singleton for every }n,
\qquad
 \#(A\cap C)=m+1,
\]
every $A\cup C_n$ is full, but $\Chat\setminus(A\cup C)$ has exactly
$m+1$ components.  Thus the limiting carrier has exactly $m$ bounded
complementary components.
\end{proposition}

\begin{proof}
Take
\[
 A=[0,m]\times\{0\}.
\]
For $n\ge1$, let $C_n$ be the polygonal arc whose ordered vertices are
\[
 (0,0),\
 \left(\tfrac12,1\right),\ (1,n^{-1}),\
 \left(\tfrac32,1\right),\ (2,n^{-1}),\ \dots,\
 \left(m-\tfrac12,1\right),\ (m,n^{-1}).
\]
The first coordinate increases strictly along the vertex list, so $C_n$ is
a Jordan arc.  It lies in the closed upper half-plane and meets the
horizontal atom $A$ only at $(0,0)$.  Both $A$ and $C_n$ are full, and their
intersection is connected; hence $A\cup C_n$ is full, either by
Janiszewski's theorem or by
Corollary~\ref{cor:full-connector-traces}.

Let $C$ be the polygonal arc obtained by replacing every vertex
$(j,n^{-1})$ by $(j,0)$.  The vertex convergence, with the same finite
linear parametrization, gives $C_n\to C$ in the Hausdorff metric and hence
$A\cup C_n\to A\cup C$.  The limit arc is still $x$-monotone and
\[
 A\cap C=\{(0,0),(1,0),\dots,(m,0)\},
\]
which has $m+1$ components.  Apply
Corollary~\ref{cor:full-connector-traces} with connector $C$ and atom $A$:
\[
 \#\pi_0\bigl(\Chat\setminus(A\cup C)\bigr)
 =1+\bigl((m+1)-1\bigr)=m+1.
\]
The component containing infinity is the unique unbounded component, so
the other $m$ components are bounded.
\end{proof}

\begin{figure}[htbp]
\centering
\begin{tikzpicture}[scale=1.05, line cap=round, line join=round]
  \begin{scope}[xshift=-3.8cm]
    \draw[line width=2.4pt,gray!70] (0,0)--(3,0);
    \draw[very thick,red!70!black]
      (0,0)--(.75,1.35)--(1.5,.20)--(2.25,1.35)--(3,.20);
    \fill (0,0) circle (1.5pt);
    \node[below] at (1.5,-.12) {$A$};
    \node[red!70!black,above] at (1.5,1.42) {$C_n$};
    \node[align=center,text width=4.4cm,font=\small] at (1.5,-1.08)
      {one attachment component\\no bounded complementary component};
  \end{scope}
  \begin{scope}[xshift=1.2cm]
    \fill[blue!12] (0,0)--(.75,1.35)--(1.5,0)--cycle;
    \fill[blue!12] (1.5,0)--(2.25,1.35)--(3,0)--cycle;
    \draw[line width=2.4pt,gray!70] (0,0)--(3,0);
    \draw[very thick,red!70!black]
      (0,0)--(.75,1.35)--(1.5,0)--(2.25,1.35)--(3,0);
    \foreach \x in {0,1.5,3} \fill (\x,0) circle (1.5pt);
    \node[below] at (1.5,-.12) {$A$};
    \node[red!70!black,above] at (1.5,1.42) {$C$};
    \node[align=center,text width=4.4cm,font=\small] at (1.5,-1.08)
      {three attachment components\\two bounded complementary components};
  \end{scope}
  \draw[-{Stealth[length=2.6mm]},thick]
    (-.15,.72)--(.65,.72) node[midway,above,font=\small] {Hausdorff limit};
\end{tikzpicture}
\caption{Trace multiplication in Proposition~\ref{prop:trace-birth}, shown
for $m=2$.  Every finite connector $C_n$ meets the atom once, so
$A\cup C_n$ is full.  In the limit, two near-contacts close and the trace
$A\cap C$ has three components.  Formula~\eqref{eq:G-full-budget} detects
the two shaded bounded complementary components exactly.}
\label{fig:trace-birth}
\end{figure}

\subsection{A safe limiting regime}
\label{subsec:G-nested-limit}

The preceding defect is specific to uncontrolled Hausdorff convergence.
There is a useful inverse-limit regime in which neither connector holes nor
trace components can be born.

\begin{theorem}[Nested-limit stability]
\label{thm:nested-limit-stability}
Let $s\ge1$, let $A_1,\dots,A_s\subset\C$ be pairwise disjoint full
continua, and let
\[
 C_1\supset C_2\supset\cdots
\]
be a decreasing sequence of nonempty full continua.  Suppose that
$C_n\cap A_i$ is nonempty and connected for every $n$ and $i$.  Then
\[
 C_\infty:=\bigcap_{n\ge1}C_n
\]
is a full continuum, every trace $C_\infty\cap A_i$ is nonempty and
connected, and
\[
 H_\infty:=C_\infty\cup\bigcup_{i=1}^{s}A_i
\]
is a full continuum.

If, in addition, the data lie in the setup of
Corollary~\ref{cor:continuum-connector}, every $C_n\Subset U$ avoids the
fixed obstacle set, and every $C_n$ satisfies that corollary's
relative-collar clearance conditions, then $H_\infty$
inherits obstacle avoidance and every protected prime-end chain survives
with the same injectivity and cyclic-order conclusions.
\end{theorem}

\begin{proof}
The nested-intersection theorem for continua makes $C_\infty$ nonempty,
compact, and connected.  Moreover
\[
 \Chat\setminus C_\infty
 =
 \bigcup_{n\ge1}(\Chat\setminus C_n).
\]
The sets on the right are connected and increasing, and all contain
$\infty$; their union is connected.  Hence $C_\infty$ is full.

For fixed $i$, the continua $C_n\cap A_i$ are nonempty, compact, connected,
and decreasing.  Their intersection is therefore nonempty and connected,
and distributivity of intersection gives
\[
 \bigcap_n(C_n\cap A_i)=C_\infty\cap A_i.
\]
Corollary~\ref{cor:full-connector-traces} now makes $H_\infty$ full.
Avoidance and relative-collar clearance pass to the subset $C_\infty$;
the final assertion follows from
Corollary~\ref{cor:continuum-connector}.
\end{proof}

The theorem suggests a concrete route for an infinite construction:
produce a \emph{decreasing} system of full connector envelopes with
connected, nonempty traces, rather than take an arbitrary subsequential
Hausdorff limit of finite trees.  It is a sufficient architecture, not a
claim that the dynamical program already supplies such envelopes.

\begin{warning}[The compactness gate]
\label{warn:compactness-gate}
Fullness and singleton attachment traces are not closed under Hausdorff
limits.  Proposition~\ref{prop:trace-birth} isolates the trace-birth
failure, while
\[
C_n=\{e^{it}:1/n\le t\le2\pi\}
\]
are closed arcs converging in Hausdorff distance to the separating Jordan
curve $S^1$, isolating the independent cycle-birth failure in the connector
itself.  Therefore a
finite protected-carrier theorem cannot be passed to an infinite or virtual
horizon by compactness alone.  At minimum, one must prove that the limiting
connector remains nonseparating and that each limiting attachment trace
remains nonempty and connected.  Under exactly those two conclusions,
Corollary~\ref{cor:full-connector-traces} restores fullness of the limiting
carrier.
\end{warning}

The warning is deliberately a no-go statement rather than a promoted
dynamical claim.  It sharpens the finite/infinite boundary of this paper:
the finite certificate solves routing, but a future limiting construction
must control both connector cohomology and trace multiplication.  Neither
control follows from Hausdorff precompactness.  Theorem
\ref{thm:nested-limit-stability} supplies one sufficient replacement:
decreasing full envelopes with nonempty connected traces.
